% Generated draft: requires mathematical and source review.
\documentclass[11pt]{article}
\usepackage[margin=25mm]{geometry}
\usepackage{fontspec}
\setmainfont{Latin Modern Roman}
\usepackage{amsmath,amssymb,mathtools,unicode-math}
\setmathfont{Latin Modern Math}
\usepackage{enumitem,needspace,xurl,hyperref}
\setlength{\parindent}{0pt}
\setlength{\parskip}{5pt}
\setlength{\emergencystretch}{3em}
\newcommand{\N}{\mathbb{N}}
\newcommand{\Z}{\mathbb{Z}}
\newcommand{\Q}{\mathbb{Q}}
\newcommand{\R}{\mathbb{R}}
\newcommand{\C}{\mathbb{C}}
\newcommand{\F}{\mathbb{F}}
\newcommand{\A}{\mathbb{A}}
\newcommand{\PP}{\mathbb{P}}
\newcommand{\E}{\mathbb{E}}
\newcommand{\eps}{\varepsilon}
\newcommand{\dd}{\,\mathrm{d}}
\begin{document}
\section*{UnsolvedMath ID: 9700002}
\textbf{Source problem code:} catalogue entry
\subsection{Definitions and mathematical statement}
The source asks for a simple network-and-demand model with a marginal satisfiability proportion $r(t)$ that can be analyzed explicitly and exhibits a proposed percolation-fragmentation congestion transition. The control parameter $t$, the marginal event defining $r(t)$, and the proposed transition are not defined in the supplied record.

A model must specify the network, demand law, capacities, and feasibility rules. In a flow interpretation these include which source-target demands may be served simultaneously, how paths or flows consume edge capacities, whether traffic can split, and what optimization objective or adaptive capacity-growth rule is used. The source offers optimization or adaptation as alternatives, not a requirement that both mechanisms be present.

The requested construction must also define the probability or limiting proportion meant by $r(t)$ and the observable that distinguishes percolation from fragmentation. A phase transition requires an appropriate infinite-volume or system-size limit, a specified order parameter, and a critical-parameter criterion; these data are not supplied. In particular $r(t)$ is not silently identified with the total fraction of all demand served, which need not be the intended marginal quantity.

Thus the task is to construct an analytically tractable model realizing the source's intended capacity/satisfiability phenomenon. The missing transition description is essential to verifying a candidate. No generic random-graph percolation transition or empirical traffic model is substituted for that requested property.
\subsection{Short English statement}
Build a simple congestion model where the intended marginal satisfiability function is analytically computable and has the specified percolation-fragmentation transition under either optimization or adaptive capacity growth. The marginal definition and transition criterion are missing.

The supplied summary reports no verified model of this exact kind and distinguishes unrelated analytic percolation results and empirical traffic studies. That is a search-based OPEN-TRIAGE assessment, not proof that no construction exists. The maintained Aldous source's current primary-literature audit remains required.
\subsection{Sources}
\begin{itemize}
\item Local UnsolvedMath registry, permanent ID 9700002, source code catalogue entry, ``Analytic toy model for a percolation-fragmentation congestion transition,'' proposed by David Aldous. Dataset curation and metadata: Ulam AI and UnsolvedMath Contributors, CC BY 4.0.
\item David Aldous, \emph{Open problems}, item \texttt{congestion.html}, index \url{https://www.stat.berkeley.edu/~aldous/Research/OP/index.html}; HTML-text extraction with linked primary-PDF text, accessed 2026-07-29. No exact linked PDF citation is supplied.
\item Provenance says the maintained index has 23 active problem pages and delegated explicit subproblems were extracted separately while compound bullets were retained. This is supplied provenance, not a new inspection claimed here.
\item The research summary mentions general percolation/fragmentation and empirical Seoul-traffic context without exact references resolving the stated model-construction problem.
\item Current primary-literature status audit and redistribution terms NEEDS\_REVIEW for the publicly accessible author/publisher copy, separately from the dataset licence.
\end{itemize}
\end{document}
